% Figure 11.3: charged-particle scattering with and without vacuum polarization.
% Source: physical PDF p. 505 (printed p. 482).
\begingroup
\renewcommand{\thefigure}{11.\arabic{figure}}
\setcounter{figure}{2}
\begin{figure}[tbp]
\centering
\begin{tikzpicture}[
  charged/.style={line width=0.65pt,
    postaction={decorate},decoration={markings,
      mark=at position 0.58 with {\arrow{Stealth[length=4.4pt,width=3.3pt]}}}},
  photon/.style={line width=0.6pt,decorate,
    decoration={snake,amplitude=1.0pt,segment length=5pt}},
  every node/.style={font=\small}
]
  \begin{scope}[xshift=-3.25cm]
    \coordinate (L) at (-1.45,0);
    \coordinate (R) at (1.45,0);
    \draw[charged] (-2.45,-1.65) -- (L);
    \draw[charged] (L) -- (-2.45,1.65);
    \draw[photon] (L) -- (R);
    \draw[charged] (2.45,-1.65) -- (R);
    \draw[charged] (R) -- (2.45,1.65);
    \node[left] at (-2.25,-1.40) {$1$};
    \node[left] at (-2.28,1.42) {$2$};
    \node[right] at (2.26,-1.40) {$1'$};
    \node[right] at (2.27,1.42) {$2'$};
    \node at (0,-2.30) {(a)};
  \end{scope}
  \begin{scope}[xshift=3.25cm]
    \coordinate (L) at (-1.55,0);
    \coordinate (R) at (1.55,0);
    \draw[charged] (-2.55,-1.65) -- (L);
    \draw[charged] (L) -- (-2.55,1.65);
    \draw[photon] (L) -- (-0.67,0);
    \draw[charged] (-0.67,0) arc[start angle=180,end angle=0,x radius=0.67cm,y radius=0.82cm];
    \draw[charged] (0.67,0) arc[start angle=0,end angle=-180,x radius=0.67cm,y radius=0.82cm];
    \draw[photon] (0.67,0) -- (R);
    \draw[charged] (2.55,-1.65) -- (R);
    \draw[charged] (R) -- (2.55,1.65);
    \node[left] at (-2.35,-1.40) {$1$};
    \node[left] at (-2.38,1.42) {$2$};
    \node[right] at (2.36,-1.40) {$1'$};
    \node[right] at (2.37,1.42) {$2'$};
    \node at (0,-2.30) {(b)};
  \end{scope}
\end{tikzpicture}
\caption{Two diagrams for the scattering of charged particles. Here lines carrying arrows are charged particles; wavy lines are photons. Diagram (b) represents the lowest-order vacuum polarization correction to the tree approximation graph (a).}
\label{fig:11.3}
\end{figure}
\endgroup
